% Part: counterfactuals
% Chapter: minimal-change-semantics
% Section: introduction

\documentclass[../../../include/open-logic-section]{subfiles}

\begin{document}

\olfileid{cnt}{min}{int}

\olsection{પરિચય}

સ્ટાલનેકર અને લૂઈસે ``જો દીવાસળી ઘસાત, તો તે સળગી ઊઠત'' જેવા
પ્રતિતથ્યાત્મક શરતી વિધાનોનું વિશ્લેષણ સૂચવ્યું. તેમનાં વિશ્લેષણો
આવાં વાક્યોની સત્ય થવાની શરતો યોગ્ય રીતે કેવી સમજવી તેની દરખાસ્તો
હતી. બંનેની પાછળનો વિચાર આ છે: પ્રતિતથ્યાત્મક શરતી વિધાન સત્ય છે
કે નહિ તે તપાસવા એવા શક્ય જગતો ધ્યાનમાં લેવાં પડે જે પૂર્વપક્ષને
સત્ય કરવા માટે વાસ્તવિક જગતથી ન્યૂનતમ પ્રમાણમાં જુદાં હોય. આ શક્ય
જગતોમાં પરિણામ સત્ય હોય, તો પ્રતિતથ્યાત્મક વિધાન સત્ય છે. દાખલા
તરીકે, ધારો કે મારા હાથમાં એક દીવાસળી અને દીવાસળીઓનું કાગળનું પૅકેટ
છે. વાસ્તવિક જગતમાં હું માત્ર તેમને જોઉં છું અને દીવાસળી ઘસું તો
શું થાય એ વિચારું છું. હું દીવાસળી ઘસું એવી, વાસ્તવિક જગતથી
ન્યૂનતમ ફેરફારવાળી સ્થિતિ એ છે જેમાં હું પગલું ભરવાનું નક્કી કરું
અને દીવાસળી ઘસું. આ ફેરફાર ન્યૂનતમ છે, કારણ કે બીજું કંઈ બદલાતું
નથી: હું સાથે હવામાં કૂદતો નથી; દીવાસળી ઘસતાં મારા વાળમાં પણ આગ
લાગતી નથી; મારી આંગળીઓની બધી તાકાત અચાનક જતી રહેતી નથી; અને
સુપર સોકર પાણીની પિસ્તોલના ઓચિંતા હુમલાથી હું એ જ સમયે ભીંજાઈ
જતો નથી, વગેરે. એ વૈકલ્પિક શક્યતામાં દીવાસળી સળગે છે. તેથી
``જો હું દીવાસળી ઘસત, તો તે સળગી ઊઠત'' એ સત્ય છે.

આ સહજ વિશ્લેષણને પ્રતિતથ્યાત્મક વિધાનોના તર્કશાસ્ત્ર માટેના ઔપચારિક
અર્થવિચાર સાથે જોડી શકાય છે. લૂઈસે પ્રતિતથ્યાત્મક સંયોજક માટે
``$\boxright$'' સંકેત રજૂ કર્યો અને સ્ટાલનેકરે~``$>$'' સંકેત વાપર્યો.
આપણે~$\cif$ વાપરીશું અને તેને વિધાનાત્મક તર્કશાસ્ત્રમાં દ્વિસ્થાની
સંયોજક તરીકે ઉમેરીશું. એટલે $!A \lif
!B$ સ્વરૂપનાં !!{formula}s ઉપરાંત~$!A \cif !B$ સ્વરૂપનાં
!!{formula}s પણ મળે છે. મોડલ તર્કશાસ્ત્રના સંબંધાત્મક અર્થવિચારની
જેમ, આ ઔપચારિક અર્થવિચાર એવા નિદર્શો પર આધારિત છે જેમાં
!!{formula}sનું મૂલ્યાંકન જગતો પર થાય છે. $\mSat{M}{!A \cif !B}[w]$
વ્યાખ્યાયિત કરતી સંતોષ-શરત કેટલાક (અન્ય) જગતો~$w'$ માટેના
$\mSat{M}{!A}[w']$ અને $\mSat{M}{!B}[w']$ વડે અપાય છે. કયા~$w'$?
સહજ રીતે કહીએ તો, $\mSat{M}{!A}[w']$ હોય એવાં~$w$ને સૌથી નજીકનાં
જગતો. આ માટે નિદર્શમાં ``નિકટતા''નો સંબંધ પણ સામેલ કરવો પડે છે.

લૂઈસે પ્રતિતથ્યાત્મક પરિસ્થિતિઓને આલેખમાં રજૂ કરવાની એક સમજણભરી
રીત રજૂ કરી. દરેક શક્ય જગત તેના આસપાસનાં અન્ય જગતોને સમાવતા
એકબીજાની અંદર આવેલા ગોળાઓના સમૂહના કેન્દ્રમાં હોય છે---આ ગોળાઓને
આપણે સમકેન્દ્રી વર્તુળો તરીકે દોરીએ છીએ. બે ગોળાઓ વચ્ચેના જગતો
કેન્દ્રના જગતથી એકસરખા નજીક હોય છે; અંદરના ગોળામાં આવેલાં જગતો
વધુ નજીક અને બહારના ગોળામાં આવેલાં જગતો વધુ દૂર હોય છે.
\begin{center}
\begin{tikzpicture}[scale=.7]\small
  \spheresystem{5}
  \draw (0,0) node {$w$};
  \propositionintersect{0}{4}{40}{3.5}
  {\tikzset{proposition/.append={smooth,tension=1.4}}
  \proposition[shift={(1.6,-1)}]{90}{1}{30}{4}}
  \path (1.5,3) node[below] {$\formula{B}$};
  \path (3,0) node {$\formula{A}$};
\end{tikzpicture}
\end{center}
સૌથી નજીકનાં $!A$-જગતો એવાં જગતો~$w'$ છે જ્યાં~$!A$ સંતોષાય છે
અને જે કેન્દ્રના જગત~$w$ની આસપાસના સૌથી નાના ગોળામાં (ભૂખરા
વિસ્તારમાં) આવેલાં છે. સહજ રીતે, સૌથી નજીકનાં બધાં $!A$-જગતોમાં
$!B$ સત્ય હોય, તો~$w$ પર $!A \cif !B$ સંતોષાય છે.

\end{document}
